\section{Метризация счётного произведения метрических пространств.}
\begin{example}
    Пусть $(X, \rho)$ -- метрическое пространство.
    Рассмотрим $X^{\N}$ -- пространство последовательностей со значениями в $X$ с топологией Тихонова.
    Покажем метризуемость этой топологии.
    Поскольку базой топологии в $X$ являются открытые шары, то предбазой в $X^{\N}$ являются множества $V(x, n, \epsilon) = \pi^{-1}_n(O_\epsilon(x(n)))$.
    Обозначим семейство таких множеств за $\sigma_T\colon$
    \[
        \sigma_T = \left\{V(x, n, \epsilon) \mid x \in X^{\N}, \ \epsilon > 0, \ n \in \N \right\}.
    \]

    Что означает, что $x \in G \in \tau_T$?
    Это означает, что существуют $\{x_i\}_{i = 1}^{m} \subset X^{\N}$, $\{n_i\}_{i = 1}^{m}$ и $\{\epsilon_i\}_{i = 1}^{m}$ такие, что
    \[
        G \supset \bigcap_{i = 1}^{m} V(x_i, n_i, \epsilon_i) \ni x.
    \]
    Поскольку мы можем взять $\epsilon = \min \epsilon_i - \rho(x(n_i), x_i(n_i))$, то
    \[
        G \supset \bigcap_{i = 1}^{m} V(x_i, n_i, \epsilon) \ni x.
    \]

    Введём метрику поточечной сходимости в $X^{\N}$.
    Заметим, что топология $(X, \rho)$ эквивалентна $(X, d)$ где $d$ определяется как
    \[
        d = \frac{\rho}{1 + \rho}.
    \]
    Все свойства метрики, кроме неравенства треугольника -- очевидны.
    Заметим, что в силу того, что $\rho \mapsto 1 - \frac{1}{1 + \rho}$ -- строго возрастает, если $\rho(a, b) \leq \rho(a, c) + \rho(b, c)$, то
    \begin{multline*}
        d(a, b) = \frac{\rho(a, b)}{1 + \rho(a, b)} \leq \frac{\rho(a, c) + \rho(c, b)}{1 + \rho(a, c) + \rho(b, c)} = \\ = \frac{\rho(a, c)}{1 + \rho(a, c) + \rho(b, c)} + \frac{\rho(b, c)}{1 + \rho(a, c) + \rho(b, c)} \leq \\ \leq \frac{\rho(a, c)}{1 + \rho(a, c)} + \frac{\rho(b, c)}{1 + \rho(b, c)} = d(a, c) + d(b, c).
    \end{multline*}
    Сходимость по $d$ совпадает со сходимостью по $\rho$, а следовательно и топологии совпадают. % Или неочевидно?

    Рассмотрим для $x, y \in X^{\N}$ метрику $d_T$, определяемую как
    \[
        d_T(x, y) = \sup_{n \in \N} \frac{d(x(n), y(n))}{n}.
    \]

    \begin{proposition}
        Если $\{x_m\} \subset X^{\N}$ и $x \in X^{\N}$, то
        \[
            x_m \xrightarrow{\tau_T} x \Longleftrightarrow x_m \xrightarrow{d_T} x.
        \]
    \end{proposition}
    \begin{proof}
        Пусть $x_m \xrightarrow{\tau_T} x$.
        Заметим, что $\forall \epsilon > 0 \exists N(\epsilon)$ такой, что $\frac{1}{N(\epsilon)} \leq \epsilon$.
        Рассмотрим
        \[
            U(x) = \bigcap_{i = 1}^{N(\epsilon)} V(x, n, \epsilon) \in \tau_T.
        \]
        Начиная с некоторого $M(\epsilon)$ все точки последовательности попадут в $U(x)$.
        А значит $\forall n \in \{1, \ldots, N(\epsilon)\}$ и $m \geq M(\epsilon)$ верно, что
        \[
            d(x(n), x_m(n)) \leq \epsilon.
        \]
        Но, тогда
        \[
            d_T(x, y) \leq \max\{\epsilon, \frac{1}{N(\epsilon)}\} \leq \epsilon.
        \]
        А значит $x_m \xrightarrow{d_T} x$.

        Пусть $x_m \xrightarrow{d_T} x$.
        Рассмотрим $U(x) \in \tau_T$.
        Существуют $\{n_i\}_{i = 1}^{N}$ и $\epsilon > 0$ такие, что
        \[
            U(x) \supset \bigcap_{i = 1}^{N} V(x, n_i, \epsilon) \ni x.
        \]
        Так как начиная с некоторого $L(\epsilon)$ метрика $d_T(x, x_m) \leq \epsilon$, то, в частности, для всякого $n_i$:
        \[
            d(x(n_i), x_m(n_i)) \leq n_i\epsilon,
        \]
        то нам достаточно выбрать максимальное из $L\left(\frac{\epsilon}{n_i}\right)$ и получить, что начиная с некоторого номера все члены последовательности $x_m$ будут лежать в $\bigcap_{i = 1}^{N} V(x, n_i, \epsilon)$.
    \end{proof}

    \begin{proposition}
        В пространстве $X^{\N}$ справедливо следующее:
        \[
            \tau_T = \tau_{d_T}.
        \]
    \end{proposition}
    \begin{proof}
        Достаточно показать, что $\sigma_T \subset \tau_{d_T}$, а, следовательно, и $\tau_{T} \subset \tau_{d_T}$.
        Аналогично, если мы покажем, что база метрической топологии -- открытые шары -- содержится в $\tau_{T}$, то и $\tau_{d_T} \subset \tau_T$.

        Рассмотрим $V(x, n, \epsilon)$ и $y \in V$.
        Если $\delta = \epsilon - d(x(n), y(n))$, то заметим, что $V(x, n, \epsilon) \supset V(y, n, \delta)$.
        Покажем, что существует $r > 0$ такое, что $O_r(y) \subset V(y, n, \delta)$.
        Заметим, что $\frac{1}{n}d(y(n), z(n)) \leq d_T(y, z)$.
        Следовательно, нам достаточно взять $r < \frac{\delta}{n}$, чтобы попасть в $V(y, n, \delta)$.
        А значит $V(x, n, \epsilon) \in \tau_{d_T}$, поскольку всякая его точка является внутренней.

        Пусть есть некоторый шар $O_R(x)$ в $(X^{\N}, d_T)$ и он содержит точку $y$.
        Тогда $O_R(x) \supset O_r(y)$, где $r = R - d_T(x, y)$.

        Нам достаточно получить, что некоторый элемент из базы $\bigcap_{n = 1}^{N} V(y, n, \epsilon)$ содержится в $O_r(y)$ при некоторых $N$ и $\epsilon$.
        Если $z \in \bigcap_{n = 1}^{N} V(y, n, \epsilon)$, то $\forall n \in \overline{1, N}$ верно, что $d(y(n), z(n)) < \epsilon$, то $d_T(y, z) < \max\{\epsilon, \frac{1}{N}\}$.
        Нам достаточно подобрать такие $\epsilon$ и $N$, чтобы $\epsilon < r$ и $\frac{1}{N} < r$.
        Пусть $\epsilon = \frac{r}{2}$ и $N = \frac{1}{r} + 1$.
        Это и завершает доказательство.
    \end{proof}
    \begin{notet}
        Заметим, что мы показали эквивалентность метрической топологии $(X, \rho)$ и $(X, d)$.
        Тогда мы можем рассматривать предбазу этой топологии именно в контексте метрики $d$ -- что несколько удобнее использования $\rho$.
    \end{notet}
    
\end{example}
\section{Декартово произведение семейства топологических пространств.}
\begin{definition}
    Пусть $T \neq \emptyset$ и $\forall t \in T$ задано некоторое $(X_t, \tau_t) \in \cat{Top}$.
    Определим $Y = \prod_{t \in T} X_t$ как
    \[
        Y = \left\{x\colon T \to \amalg_{t \in T} X_t \mid x(t)\in X_t \right\}.
    \]
    Определим канонические проекции как
    \[
        \pi_t\colon Y \to X_t, x \mapsto x(t).
    \]
    Рассмотрим семейство $\sigma_T$, определяемое как
    \[
        \sigma_T = \left\{\pi^{-1}_t(G_t) \mid \forall t \in T, \forall G_t \in \tau_t\right\}.
    \]
    Она является предбазой топологии Тихонова в $Y$ (очевидно, что она удовлетворяет критерию предбазы).

    Отображение $\pi_t$ в топологии, порождённой данной предбазой являются непрерывными.
\end{definition}
\begin{proposition}
    Топология Тихонова -- слабейшая топология, относительно которой все канонические проекции являются непрерывными.
\end{proposition}
\begin{proof}
    Пусть $\hat{\tau}$ -- слабейшая топология на $Y$ такая, что все канонические проекции непрерывны.
    Тогда множества вида $\pi^{-1}_t(G_t)$ для всякого $G_t \in \tau_t$ и $t \in T$ являются открытыми в ней.
    Следовательно, $\sigma_T \subset \hat{\tau}$ и $\tau_T \subset \hat{\tau}$.
    Поскольку в $\tau_T$ все канонические проекции непрерывны, то $\tau_T = \hat{\tau}$.
\end{proof}
\begin{notet}
    Я не понял зачем Конст делает то, что делает.
    Ему по кайфу, видимо.
    Я перетехаю этот кусок.
    Ну и про поточечную сходимость аналогичное доказательство (абсолютно аналогичное тому, что было в декартовом произведении одинаковых пространств).
\end{notet}
\begin{note}
    Утверждение про метризуемость счётного произведения метрических пространств можно обобщить на произведение различных пространств.
    Абсолютно аналогично.
\end{note}

\chapter{Компактность.}
\begin{definition}
    Пусть $P \subset 2^{X}$ и $(X, \tau) \in \cat{Top}$.

    Если для $K \subset X$ верно, что $K \subset \bigcup_{V \in P} V$, то $P$ -- покрытие $K$.

    Если $P \subset \tau$, к тому же, то будем называть $P$ -- открытым покрытием $K$.
\end{definition}
\begin{definition}
    Пусть $(X, \tau) \in \cat{Top}$.
    Будем говорить что множество $K \subset X$ является компактным, если из всякого его открытого покрытия можно выделить конечное подпокрытие.%\footnote{Традиционно, компактность не определяется для множества в топологическом пространстве -- она даётся как характеристика самого пространства и далее, множество называется компактом, если оно будет компактно в индуцированной топологии. Впрочем, Роман Викторович существование топологии подпространства упорно игнорирует.}
\end{definition}
\begin{definition}
    Пусть $(X, \tau) \in \cat{Top}$. 
    Будем говорить что множество $K \subset X$ является секвенциально компактным, если из всякой последовательности его элементов можно выделить сходящуюся подпоследовательность.%\footnote{А вот это свойство топологическим не является, и через топологию подпространства в общем случае не может быть определено. Впрочем, оно нужно только для метрических пространств -- а там можно определить и так.}
\end{definition}
Сформулируем теорему Тихонова.
\begin{theorem}[Тихонов]
    Если есть некоторое семейство топологических пространств $\{X_t\}_{t \in T}$, где $T$ -- непустое индексное множество,
    и при всяком $t \in T$ множество $K_t$ является компактом, то $K = \prod_{t \in T} K_t$ является компактом в топологии Тихонова в $Y = \prod_{t \in T} X_t$.
\end{theorem}
\begin{example}
    Гильбертов куб (то есть $[0, 1]^{[0, 1]}$) компактен в силу теоремы Тихонова.
\end{example}
\begin{example}
    Гильбертов куб не является секвенциально компактным.

    Пусть $x_n\colon [0, 1] \to [0, 1]$ и $x_n(t)$ -- $n$-ая цифра числа $t$ в двоичной записи.
    Предположим, что существует сходящаяся подпоследовательность $\{x_{n_k}\}$.
    Рассмотрим $t^* \in [0, 1]$ -- такое число, что $n_{2k}$-ая цифра этого числа есть единица, а остальные -- нули.
    Тогда $x_{n_k}(t^*)$ -- расходится, поскольку имеет два частичных предела -- по чётным членам подпоследовательности -- 1, а по нечётным -- 0.
\end{example}
\section{Некоторые сведения из теории множеств.}
\begin{definition}
    Будем говорить, что пара $(X, \preceq)$, где $X$ -- множество, а $\preceq$ -- бинарное отношение на нём, является частично упорядоченным множеством, если
    \begin{enumerate}
        \item $\forall x, y \in X x \preceq y \wedge y \preceq x \Rightarrow x = y$;
        \item $\forall x, y, z \in X\colon x \preceq y, \ y \preceq z  \Rightarrow x \preceq z$;
        \item $\forall x \in X \hookrightarrow x \preceq x$.
    \end{enumerate}
\end{definition}
\begin{note}
    Любое семейство $\mathcal{F} \subset 2^X$ превращается в частично упорядоченное, если в качестве порядка взять вложение.
\end{note}
\begin{definition}
    Цепью в частично упорядоченном множестве $(X, \preceq)$ называется линейно-упорядоченное подмножество (то есть такое подмножество, что любые два его элемента сравнимы).
\end{definition}
\begin{lemma}[Цорн]
    Если для всякой цепи $C$ в частично упорядоченном множестве $(X, \preceq)$ существует мажоранта (то есть элемент $p$, такой, что $\forall x \in C \hookrightarrow x \preceq p$), то в $X$ существует максимальный элемент (такой элемент $p$, что $\nexists x \neq p \in X\colon p \preceq x$).
\end{lemma}
\begin{theorem}[Хаусдорф]
    Всякое частично упорядоченное множество имеет максимальную (по вложению) цепь.
\end{theorem}
\begin{proposition}
    Лемма Цорна и теорема Хаусдорфа эквивалентны.
\end{proposition}
\begin{proof}
    Пусть выполнена теорема Хаусдорфа.
    Рассмотрим $(X, \preceq)$ -- частично упорядоченное множество такое, что у всякой цепи существует мажоранта.
    Пусть $C_*$ -- максимальная цепь в $X$.
    Существует некоторый элемент $p \in X\colon \forall x \in C_* \hookrightarrow x \preceq p$.

    Покажем, что $p$ -- максимальный элемент.
    Пусть это не так.
    Тогда существует некоторый элемент $q \in X\colon p \preceq q$ и $p \neq q$.
    В силу транзитивности $\forall x \in C_* \hookrightarrow x \preceq q$.
    Тогда $C_* \cup \{q\}$ -- цепь.
    Но $C_*$ -- максимальна по включению, а значит $q \in C_*$.
    Тогда $q \preceq p$ и $p \preceq q$ -- противоречие и $p$ -- максимальный элемент.

    Пусть теперь выполнена лемма Цорна.
    Рассмотрим произвольное частично упорядоченное множество $(X, \preceq)$ и рассмотрим его пространство цепей $(Y, \subseteq)$.
    Пространство цепей допускает естественное упорядочивание по вложению и, очевидно, что для произвольной цепи из цепей объединение цепей является мажорантой цепи из цепей -- а значит выполнены условия леммы Цорна и существует максимальный по включению элемент в пространстве цепей.
    Он и будет максимальной цепью. % допилить
\end{proof}

\section{Доказательство теоремы Хаусдорфа (б/д, как я понимаю).}
\begin{notet}
    За тех этой секции большое спасибо Владу Самсонову.
\end{notet}

Доказательство будет устроено так:
\begin{itemize}
    \item Если докажем, что существует цепь $C$, для которой $C^{*}=\varnothing$, то победим.
    \item Если докажем лемму о неподвижном множестве, то найдём элемент с пустым $C^{*}$.
    \item Если покажем, что минимальная башня $\mathcal T^{*}$ является цепью, то докажем лемму о неподвижном множестве.
    \item Если покажем, что специально построенное $\Gamma\subset\mathcal T^{*}$ является башней, то докажем, что $\mathcal T^{*}$ --- цепь.
    \item Если покажем, что специально построенное $M(A)=\mathcal T^{*}$, то докажем, что $\Gamma$ --- башня.
    \item Наконец, покажем, что $M(A)$ --- башня, из чего последует $M(A)=\mathcal T^{*}$.
\end{itemize}

\subsection{}*{3. Шаг 1: сведём всё к цепи с $C^{*}=\varnothing$}

Пусть $X$ --- носитель ЧУМ $(X,\le)$ и
\[
    F=\{C\subset X:\ C \ \text{--- цепь в }(X,\le)\}.
\]
Тогда $F\neq\varnothing$, поскольку $\{x\}\in F$ для любого $x\in X$.

Для каждой цепи $C\in F$ определим множество ``расширений на один элемент'':
\[
    C^{*}=\{x\in X\setminus C:\ C\cup\{x\}\in F\}.
\]
Тогда:
\[
    C\ \text{максимальна}\quad\Longleftrightarrow\quad C^{*}=\varnothing.
\]
Действительно, если $C^{*}\neq\varnothing$, то можно добавить $x\in C^{*}$ и получить большую цепь.
Если же $C^{*}=\varnothing$, то добавить нельзя ни один элемент, значит цепь максимальна.

Значит, достаточно построить $C\in F$ с $C^{*}=\varnothing$.

\subsection{}*{4. Шаг 2: строим отображение $g$ и формулируем лемму о неподвижном множестве}

\subsubsection*{4.1. Аксома выбора и построение $g$}
Используем аксиому выбора в форме: для семейства всех непустых подмножеств $X$ существует функция выбора
$f$, такая что $f(A)\in A$ для любого $A\subset X$, $A\neq\varnothing$.

Для $C\in F$ определим отображение $g:F\to F$ так:
\[
    g(C)=
    \begin{cases}
        C, & C^{*}=\varnothing,\\
        C\cup\{f(C^{*})\}, & C^{*}\neq\varnothing.
    \end{cases}
\]
Тогда для любого $C\in F$ выполнено:
\[
    C\subset g(C)\quad\text{и}\quad |g(C)\setminus C|\le 1.
\]
И главное: если $g(C)=C$, то обязательно $C^{*}=\varnothing$, т.е. $C$ максимальна.

\subsubsection*{4.2. Лемма о неподвижном множестве}
\begin{lemma}[о неподвижном множестве]
    Пусть $(F,\subset)$ --- ЧУМ и $g:F\to F$ таково, что
    \[
        \forall C\in F:\quad C\subset g(C),\qquad |g(C)\setminus C|\le 1.
    \]
    Тогда существует $C_{\infty}\in F$ такое, что $g(C_{\infty})=C_{\infty}$.
\end{lemma}

\noindent
Из этой леммы сразу следует теорема Хаусдорфа: берём $C_{\infty}$, тогда $g(C_{\infty})=C_{\infty}$
возможно только при $C_{\infty}^{*}=\varnothing$, значит $C_{\infty}$ --- максимальная цепь.

Остаётся доказать лемму о неподвижном множестве.

\subsection*{5. Шаг 3: башни и минимальная башня $\mathcal T^{*}$}

Зафиксируем некоторое $C_0\in F$ (например, $C_0=\{x_0\}$ для произвольного $x_0\in X$).

\begin{definition}[Башня над $C_0$]
    Подмножество $\mathcal T\subset F$ называется \emph{башней над $C_0$}, если:
    \begin{enumerate}
        \item $C_0\in\mathcal T$ и $\forall C\in\mathcal T:\ C_0\subset C$;
        \item если $\mathcal S\subset\mathcal T$ --- цепь в $(F,\subset)$, то $\displaystyle \bigcup_{C\in\mathcal S} C\in\mathcal T$;
        \item $\forall C\in\mathcal T:\ g(C)\in\mathcal T$.
    \end{enumerate}
\end{definition}

Рассмотрим пересечение всех башен над $C_0$:
\[
    \mathcal T^{*}=\bigcap_{\mathcal T\ \text{--- башня над }C_0}\mathcal T.
\]
Стандартной проверкой по пунктам (1)--(3) получаем, что $\mathcal T^{*}$ само является башней
(минимальной по включению среди башен над $C_0$).

\subsubsection*{Ключевая редукция}
Если мы докажем, что $\mathcal T^{*}$ является \emph{цепью} в $(F,\subset)$, то, положив
\[
    C_{\infty}=\bigcup_{C\in\mathcal T^{*}} C,
\]
получим:
\begin{itemize}
    \item $C_{\infty}\in\mathcal T^{*}$ по (2), потому что $\mathcal T^{*}$ цепь и можно брать объединение;
    \item $g(C_{\infty})\in\mathcal T^{*}$ по (3);
    \item но $C_{\infty}$ есть объединение всех элементов $\mathcal T^{*}$, поэтому для любого $D\in\mathcal T^{*}$
    имеем $D\subset C_{\infty}$, а значит и $g(C_{\infty})\subset C_{\infty}$ (так как $g(C_{\infty})\in\mathcal T^{*}$);
    \item с другой стороны, всегда $C_{\infty}\subset g(C_{\infty})$.
\end{itemize}
Следовательно, $g(C_{\infty})=C_{\infty}$, что и доказывает лемму о неподвижном множестве.

Значит, осталось: \emph{доказать, что $\mathcal T^{*}$ --- цепь.}

\subsection*{6. Шаг 4: строим $\Gamma\subset\mathcal T^{*}$ и доказываем, что это башня}

Определим
\[
    \Gamma=\bigl\{A\in\mathcal T^{*}:\ \forall B\in\mathcal T^{*}\ \ (A\subset B\ \text{или}\ B\subset A)\bigr\}.
\]
То есть $\Gamma$ --- элементы $\mathcal T^{*}$, сравнимые по включению со \emph{всеми} элементами $\mathcal T^{*}$.

Заметим: $C_0\in\Gamma$, поскольку для любого $B\in\mathcal T^{*}$ из свойства башни следует $C_0\subset B$.

\medskip

\noindent\textbf{Идея:} если $\Gamma$ --- башня, то по минимальности $\mathcal T^{*}$ получим
$\mathcal T^{*}\subset\Gamma$, а так как $\Gamma\subset\mathcal T^{*}$ по определению, то $\Gamma=\mathcal T^{*}$.
Тогда любой элемент $\mathcal T^{*}$ сравним со всеми, значит $\mathcal T^{*}$ --- цепь.

\medskip

\noindent
Проверим, что $\Gamma$ --- башня.

\begin{enumerate}
    \item $C_0\in\Gamma$ уже отмечено.
    \item Пусть $\mathcal S\subset\Gamma$ --- цепь (по включению). Положим $A=\bigcup_{D\in\mathcal S}D$.
    Тогда $A\in\mathcal T^{*}$ (так как $\mathcal T^{*}$ --- башня).
    Покажем $A\in\Gamma$. Возьмём любое $B\in\mathcal T^{*}$.
    Для каждого $D\in\mathcal S$ (а $D\in\Gamma$) имеем: $D\subset B$ или $B\subset D$.
    Если для некоторого $D$ верно $B\subset D$, то $B\subset A$. Иначе для всех $D\in\mathcal S$ верно $D\subset B$,
    тогда $A=\bigcup D\subset B$. Значит, $A$ сравнимо с $B$, то есть $A\in\Gamma$.
    \item Осталось: $\forall A\in\Gamma\ \ g(A)\in\Gamma$.
    Это самое тонкое место; для него вводится множество $M(A)$.
\end{enumerate}

\subsection*{7. Шаг 5: вводим $M(A)$ и сводим $g(A)\in\Gamma$ к равенству $M(A)=\mathcal T^{*}$}

Зафиксируем $A\in\Gamma$. Определим
\[
    M(A)=\{B\in\mathcal T^{*}:\ B\subset A\ \ \text{или}\ \ g(A)\subset B\}.
\]
Очевидно, $M(A)\subset\mathcal T^{*}$.

\medskip

\noindent
Если мы докажем, что $M(A)$ --- башня над $C_0$, то по минимальности $\mathcal T^{*}$ будет
\[
    \mathcal T^{*}\subset M(A)\subset\mathcal T^{*}\quad\Rightarrow\quad M(A)=\mathcal T^{*}.
\]
А тогда для любого $B\in\mathcal T^{*}$ верно:
\[
    B\subset A\quad\text{или}\quad g(A)\subset B.
\]
В частности, если $A\subset B$, то первая альтернатива невозможна (если $A\subsetneq B$, то $B\not\subset A$),
значит обязательно $g(A)\subset B$. То есть $g(A)$ сравнимо с любым $B\in\mathcal T^{*}$.
Следовательно, $g(A)\in\Gamma$.

Итак, \textbf{достаточно доказать, что $M(A)$ --- башня.}

\subsection*{8. Шаг 6: доказываем, что $M(A)$ --- башня, значит $M(A)=\mathcal T^{*}$}

Проверим свойства башни для $M(A)$.

\subsubsection*{(1) $C_0\in M(A)$}
Так как $A\in\mathcal T^{*}$ и $\mathcal T^{*}$ --- башня, имеем $C_0\subset A$. Значит $C_0\in M(A)$
(по первому условию $B\subset A$).

\subsubsection*{(2) Замкнутость относительно объединения цепей}
Пусть $\mathcal S\subset M(A)$ --- цепь по включению. Рассмотрим $U=\bigcup_{B\in\mathcal S}B$.
Так как $\mathcal S\subset\mathcal T^{*}$ и $\mathcal T^{*}$ --- башня, то $U\in\mathcal T^{*}$.

Если найдётся $B\in\mathcal S$ такое, что $g(A)\subset B$, то $g(A)\subset U$, значит $U\in M(A)$.
Иначе для всех $B\in\mathcal S$ имеем $B\subset A$, тогда и $U\subset A$, значит $U\in M(A)$.

\subsubsection*{(3) Замкнутость относительно $g$}
Пусть $B\in M(A)$. Нужно доказать $g(B)\in M(A)$.
Так как $B\in\mathcal T^{*}$ и $\mathcal T^{*}$ --- башня, то $g(B)\in\mathcal T^{*}$.

Рассмотрим два случая.

\medskip
\noindent\textbf{Случай 1: $g(A)\subset B$.}
Тогда $g(A)\subset B\subset g(B)$ (ведь всегда $B\subset g(B)$), значит $g(A)\subset g(B)$.
Следовательно, $g(B)\in M(A)$.

\medskip
\noindent\textbf{Случай 2: $B\subset A$.}
Если $g(B)\subset A$, то $g(B)\in M(A)$ (первое условие).
Пусть теперь $g(B)\not\subset A$. Тогда, поскольку $A\in\Gamma$ и $g(B)\in\mathcal T^{*}$,
элементы $A$ и $g(B)$ сравнимы, значит при $g(B)\not\subset A$ обязательно
\[
    A\subset g(B).
\]
Итак, имеем цепочку включений
\[
    B\subset A\subset g(B).
\]
Заметим, что по условию на $g$ множество $g(B)\setminus B$ содержит не более одного элемента.
Если бы одновременно были строгие включения $B\subsetneq A$ и $A\subsetneq g(B)$, то существовали бы
элемент из $A\setminus B$ и элемент из $g(B)\setminus A$, то есть $g(B)\setminus B$ содержало бы
как минимум два элемента --- противоречие. Значит, одно из включений является равенством:
\[
    \text{либо }A=B,\quad \text{либо }A=g(B).
\]
Если $A=g(B)$, то $g(B)\subset A$ и мы уже закончили.
Если $A=B$, то $g(B)=g(A)$, а значит $g(A)\subset g(B)$ (даже равенство), т.е. $g(B)\in M(A)$.

\medskip

Во всех случаях $g(B)\in M(A)$. Значит, $M(A)$ --- башня.

\subsubsection*{Итог}
Раз $M(A)$ --- башня и $M(A)\subset\mathcal T^{*}$, то по минимальности $\mathcal T^{*}$:
\[
    M(A)=\mathcal T^{*}.
\]
Отсюда следует $g(A)\in\Gamma$, значит $\Gamma$ --- башня, значит $\Gamma=\mathcal T^{*}$,
значит $\mathcal T^{*}$ --- цепь.

Тогда построенное $C_{\infty}=\bigcup_{C\in\mathcal T^{*}}C$ удовлетворяет $g(C_{\infty})=C_{\infty}$,
а значит $C_{\infty}^{*}=\varnothing$, т.е. $C_{\infty}$ --- максимальная цепь.

\medskip
\begin{center}
    \textbf{Теорема Хаусдорфа доказана.}
\end{center}

\section{Теорема Тихонова}
\begin{lemma}[Александер]
    Пусть $(X, \tau) \in \cat{Top}$, $\sigma$ -- предбаза топологии $\tau$ и $K \subset X\colon$ для всякого $\sigma$-покрытия $K$ имеет конечное подпокрытие.
    Тогда $K$ -- компактно в $X$.
\end{lemma}
\begin{proof}
    Предположим, что $K$ не является компактом.
    Тогда существует такое открытое покрытие $K$, что из него нельзя выделить конечное подпокрытие.
    И семейство открытых покрытий $F$, не имеющих конечного подпокрытия непусто:
    \[
        F = \left\{P \subset \tau \mid P\text{ не допускает конечное подпокрытие. }\right\}.
    \]
    Упорядочим $F$ по вложению.
    Заметим, что для всякой цепи $C \subset F$ покрытие $P_C = \bigcup_{P \in C} P$ также лежит в $F$ и, следовательно, всякая цепь имеет мажоранту.
    Согласно лемме Цорна $F$ имеет максимальный элемент $P^*$.
    Это означает, что если мы добавим в $P^*$ произвольное открытое $U \in \tau$, не лежащее в нём, то $P^* \cup \{U\} \notin F$.
    Значит $\forall U \in \tau \setminus P^*$ существует конечное подпокрытие $K$ -- некоторые $\{V_k\}_{k = 1}^{N} \subset P^*$ такие, что
    $\{U\} \cup \{V_k\}$ -- покрытие $K$.

    Вспомним, что по условию теоремы всякое покрытие элементами предбазы имеет конечное подпокрытие.
    Пусть $P_\sigma = P^* \cap \sigma$ -- элементы предбазы, лежащие в максимальном покрытии.
    Заметим, что $P_\sigma$ не может быть покрытием $K$ -- иначе у него существует конечное подпокрытие.
    Значит существует $x^* \in K$, не покрытая $P_\sigma$.
    Поскольку $P^*$ -- покрытие, то существует $V^* \in P^*$ такое, что $x^* \in V^*$.
    Следовательно, существуют элементы из предбазы $\{W_j\}_{j = 1}^{M} \subset \sigma$ такие, что $x^* \in \cap_{j = 1}^{M} W_j \subset V^*$.
    Поскольку $x^*$ не покрывается элементами $P_{\sigma}$, то всякое $W_j \notin P^*$.

    Тогда $\forall k \in \{1, \ldots, M\}$ существуют $\{V_{k, j}\}_{j = 1}^{N_k}$ такие, что
    \[
        W_k \cup \bigcup_{j = 1}^{N_k} V_{k, j} \supset K.
    \]
    А значит
    \[
        K \subset \bigcap_{k = 1}^{M} W_k \cup \bigcup_{k, j}^{M, N_k} V_{k, j} \subset V^* \cup \bigcup_{k, j}^{M, N_k} V_{k, j}.
    \]
    Но это конечное подпокрытие $K$ элементами $P^*$ -- чего быть не может.
    Следовательно, всякое открытое покрытие $K$ допускает конечное подпокрытие и $K$ -- компакт. % допилить
\end{proof}
\begin{theorem}[Тихонов]
    Если есть некоторое семейство топологических пространств $\{X_t\}_{t \in T}$, где $T$ -- непустое индексное множество,
    и при всяком $t \in T$ множество $K_t$ является компактом, то $K = \prod_{t \in T} K_t$ является компактом в топологии Тихонова в $Y = \prod_{t \in T} X_t$.
\end{theorem}
\begin{proof}
    Известно, что топология Тихонова имеет следующую предбазу $\sigma_T\colon$
    \[
        \sigma_T = \left\{\pi^{-1}_t(U) \mid \forall t \in T \ \forall U \in \tau_t, \text{ где }\tau_t \text{ -- топология на }X_t\right\},
    \]
    а под отображениями $\pi_t$ понимаются канонические проекции, то есть
    \[
        \pi_t\colon \prod_{t \in T} X_t \to X_t, \ x \mapsto x(t).
    \]
    Покажем то, что мы можем из всякого открытого покрытия $K$ элементами предбазы выбрать конечное подпокрытие.
    Для всякого $t \in T$ рассмотрим семейство множество $\beta_t\colon$
    \[
        \beta_t = \left\{G_t \mid G_t \in \tau_t, \ \pi^{-1}_t(G_t) \in \beta\right\}.
    \]
    Заметим, что существует $t_0 \in T$ такое, что $\beta_{t_0}$ является покрытием $K_{t_0}$ (поскольку в обратном случае для всякого $t \in T$ существует точка $y_t$ такая, что она не пересекается с объединением элементов $\beta_t$ и, следовательно, существует функция, которая принимает значения в $K_t$ на всяком $t$, но не покрывается элементами $\beta$ -- противоречие).

    Поскольку $\beta_{t_0}$ -- является покрытием $K_{t_0}$ и $K_{t_0}$ -- компактно, то существует конечное подпокрытие $\{G_k\}_{k = 1}^{N}$.
    Тогда и $\{V_k\} = \{\pi^{-1}_t(G_k)\}$ является конечным подпокрытием $\beta$ для $K$.

    Следовательно, в силу леммы Александера, $K$ является компактом.
\end{proof}